\documentclass[11pt,letterpaper]{article}
\usepackage[T1]{fontenc}
\usepackage[utf8]{inputenc}
\usepackage{lmodern,microtype}
\usepackage[margin=0.79in,headheight=15pt]{geometry}
\usepackage{parskip,enumitem,longtable,tabularx,booktabs,array,xcolor,fancyhdr,hyperref}
\definecolor{navy}{HTML}{17314C}
\definecolor{teal}{HTML}{087F78}
\hypersetup{colorlinks=true,linkcolor=navy,urlcolor=teal}
\setlength{\emergencystretch}{3em}
\setlist[itemize]{leftmargin=*,nosep,topsep=5pt}
\setlist[enumerate]{leftmargin=*,itemsep=3pt}
\renewcommand{\arraystretch}{1.22}
\pagestyle{fancy}\fancyhf{}
\fancyhead[L]{\small Enterprise Data Modeling}\fancyhead[R]{\small Documentation Suite}
\fancyfoot[C]{\thepage}
\newcommand{\srcbook}{\textit{The Data Model Resource Book}, revised edition, Volume 2}
\newcommand{\ResponseLine}{\par\vspace{1.05em}\noindent\rule{\linewidth}{0.4pt}\par}
\newcommand{\ResponseArea}{\ResponseLine\ResponseLine\ResponseLine}
\begin{document}
\begin{center}
{\LARGE\bfseries\color{navy} Universal Data Modeling Patterns}\par
\vspace{0.25em}{\large Conceptual, logical and physical modeling}\par
\vspace{0.3em}{\small Original engineering study material \textbullet{} 2026 edition}
\end{center}
\vspace{0.8em}
\section{Purpose and authority}
This is a reusable modeling vocabulary inspired by the book introduction and cross-industry themes. It is not a copy of its diagrams or a formal standard. Treat the model as a collection of reusable concepts that must be validated against an organization's actual semantics and requirements.


\section{Cross-industry building blocks}
\begin{tabularx}{\linewidth}{@{}p{0.23\linewidth} p{0.39\linewidth} p{0.34\linewidth}@{}}
\toprule
\textbf{Concept} & \textbf{Typical responsibility} & \textbf{Critical question} \\ 
\midrule
Party & Identity for person or organization & Are roles and names independently versioned? \\ 
Party role & A party acting in a context & Can one party have simultaneous roles? \\ 
Relationship / assignment & Association with qualifiers, role or time & Does the relation require its own history? \\ 
Product / offering & Catalog definition separated from commercial terms & Can a service be bundled with a good? \\ 
Agreement & Mutual commitment, terms and parties & Which clauses change and how are revisions recorded? \\ 
Order / order item & Requested commitment with line-level facts & When does the order become binding? \\ 
Shipment / deployment & Operational fulfillment or installation & How are partial events represented? \\ 
Work effort & Activities, actuals, planned allocations & What is measured and at what granularity? \\ 
Transaction & Immutable business posting or event & How are reversals and idempotency handled? \\ 
Time / unit & Temporal dimensions and measures & Does the unit have a defined conversion? \\ 
\bottomrule
\end{tabularx}

\section{Logical modeling discipline}
\begin{itemize}
\item Separate entity and attribute names from their physical SQL representations.
\item Specify candidate keys, stable identifiers, alternate keys and uniqueness scope.
\item Declare cardinality, optionality, generalization completeness and overlap policy.
\item Treat roles as relations or typed assignments when time-varying.
\item Separate conceptual containment from physical cascade behavior.
\item Distinguish valid time from recorded time when history is legally or operationally important.
\end{itemize}

\section{Physical-model checklist}
\begin{itemize}
\item Implement primary keys, unique constraints and foreign keys.
\item Use fixed-precision numeric values for financial amounts and explicit currencies.
\item Measure query access patterns before choosing indexes.
\item Use transactions to protect cross-row invariants; check constraints cannot represent arbitrary cross-row logic.
\item Review privacy, retention, lineage and access controls at table and column level.
\end{itemize}

\section{Mapping workflow}
\begin{enumerate}
\item Discover domain vocabulary and example records.
\item Identify concepts common across domains but not necessarily common physical schemas.
\item Choose variation points and write invariants.
\item Render ERD and UML class models with documented semantic conventions.
\item Implement and migrate a small physical schema.
\item Compare negative tests and sample queries against business expectations.
\end{enumerate}

\vfill
\hrule
{\footnotesize\textbf{Provenance.} Len Silverston, The Data Model Resource Book, Revised Edition, Volume 2: A Library of Universal Data Models by Industry Types (Wiley, 2001). Introduction, pp. 1--26, and Chapter 10, pp. 435--444; supplementary engineering interpretation.}
\end{document}
